\documentclass[11pt,a4paper]{article}

\usepackage[margin=2.4cm]{geometry}
\usepackage{amsmath,amssymb}
\usepackage{booktabs}
\usepackage{array}
\usepackage{xcolor}
\usepackage[colorlinks=true,linkcolor=blue!50!black,urlcolor=blue!55!black]{hyperref}
\usepackage{enumitem}
\usepackage{titlesec}

\titleformat{\section}{\normalfont\Large\bfseries\color{blue!45!black}}{\thesection}{1em}{}
\titleformat{\subsection}{\normalfont\large\bfseries\color{blue!30!black}}{\thesubsection}{1em}{}

\newcommand{\Rey}{\mathrm{Re}}
\newcommand{\bu}{\mathbf{u}}
\newcommand{\bv}{\mathbf{v}}
\newcommand{\bx}{\mathbf{x}}
\newcommand{\bn}{\mathbf{n}}
\newcommand{\ba}{\mathbf{a}}
\newcommand{\bd}{\mathbf{d}}
\newcommand{\bF}{\mathbf{F}}
\newcommand{\bR}{\mathbf{R}}
\newcommand{\bff}{\mathbf{f}}
\newcommand{\bt}{\mathbf{t}}
\newcommand{\bg}{\mathbf{g}}
\newcommand{\bJ}{\mathsf{J}}
\newcommand{\bA}{\mathsf{A}}
\newcommand{\bI}{\mathsf{I}}
\newcommand{\bttau}{\boldsymbol{\tau}}
\newcommand{\dd}{\,\mathrm{d}}
\newcommand{\mdot}{\dot m}

\title{\bfseries The CVFEM Navier--Stokes operators\\
\large Discretisation of the incompressible equations on trilinear hexahedra,\\
and the family of discrete operators a matrix-free solver evaluates}
\date{\today}

\begin{document}
\maketitle

\begin{abstract}
\noindent
This note fixes the discretisation and enumerates the discrete operators built on it. The
scheme is a colocated control-volume finite-element method on trilinear hexahedra: the
unknowns live at the mesh nodes, the balance is written over the dual control volumes, and
the fluxes are evaluated on sub-control surfaces cut out of each element. Pressure--velocity
coupling is by a Rhie--Chow mass-flux correction; convection is upwinded through a switch
that can be smoothed to keep the flux differentiable at stagnation.

The reason to write the operators down separately, rather than as one residual and its
derivative, is that they are not one thing. A Newton--Krylov solve evaluates at least five
of them per step, they carry different terms by design rather than by omission, and they
differ in cost by a factor of fifty. Section~\ref{sec:operators} lists them with their
domains of definition, and Section~\ref{sec:cost} with their measured throughput.
\end{abstract}

\section{Continuous problem}

On a bounded domain $\Omega \subset \mathbb{R}^3$ we solve the incompressible Navier--Stokes
equations in conservative form for velocity $\bu$ and pressure $p$, with constant density
$\rho$ and dynamic viscosity $\mu$,
\begin{align}
  \rho \frac{\partial \bu}{\partial t}
  + \nabla \cdot \left( \rho\, \bu \otimes \bu \right)
  + \nabla p - \nabla \cdot \bttau &= \rho\, \bff , \label{eq:mom} \\[2pt]
  \nabla \cdot \bu &= 0 , \label{eq:cont}
\end{align}
with the Newtonian deviatoric stress
\begin{equation}
  \bttau = \mu \left( \nabla \bu + \nabla \bu^{\mathsf T} \right) ,
  \qquad \tau_{cd} = \mu \left( \partial_d u_c + \partial_c u_d \right) .
  \label{eq:stress}
\end{equation}
The trace term $-\tfrac23 \mu (\nabla\cdot\bu)\,\bI$ is absent: it vanishes for a
divergence-free field and retaining it would make the discrete stress depend on how well
\eqref{eq:cont} is satisfied at the current iterate rather than at the solution.

Integrating \eqref{eq:mom}--\eqref{eq:cont} over a control volume $C$ with boundary
$\partial C$ and applying the divergence theorem gives the balance the discretisation
reproduces:
\begin{align}
  \frac{\dd}{\dd t} \int_C \rho\, \bu \dd V
  + \oint_{\partial C} \left[ \rho\, \bu (\bu \cdot \bn) + p\,\bn - \bttau \bn \right] \dd S
  &= \int_C \rho\, \bff \dd V , \label{eq:cvmom} \\[2pt]
  \oint_{\partial C} \rho\, \bu \cdot \bn \dd S &= 0 . \label{eq:cvcont}
\end{align}
Every term of every operator below is a quadrature of one of these two, so the discrete
system is conservative by construction: a flux leaving one control volume enters its
neighbour with the opposite sign and the same value.

\section{The dual mesh and its geometry}

\subsection{Control volumes and sub-control surfaces}

The primal mesh is a partition of $\Omega$ into trilinear hexahedra. The unknowns
$(\bu_i, p_i)$ are colocated at the $N$ mesh nodes, and the control volume $C_i$ is the
median-dual cell of node $i$: within each element containing $i$, the sub-volume bounded by
the planes through the element centroid, the face centroids and the edge midpoints. Each
hexahedron therefore contributes one of its eight sub-volumes to each of its eight nodes,
and
\begin{equation}
  V_i \;=\; \sum_{e \,\ni\, i} \frac{1}{8} \left| \det \bJ_e \right| ,
  \label{eq:nodevol}
\end{equation}
where $\bJ_e$ is the element Jacobian defined below. The $V_i$ partition $\Omega$ exactly:
$\sum_i V_i = |\Omega|$.

The interface between $C_i$ and $C_j$ inside one element is a \emph{sub-control surface}
(SCS). In a hexahedron the pairs $(i,j)$ that share such an interface are exactly the twelve
edges, so each element carries $N_{\mathrm{scs}} = 12$ interior sub-control surfaces, grouped
four to each of the three parametric directions. A sub-control surface lying on
$\partial\Omega$ is instead a \emph{boundary sub-face}; each of the six element faces carries
four of them, one per node of that face, and they close the control volumes that touch the
domain boundary.

\subsection{Isoparametric map and area vectors}

Let $\bx_a$, $a = 0,\dots,7$, be the nodes of an element and $N_a(\boldsymbol{\xi})$ the
trilinear shape functions on the reference cube $\hat\Omega = [0,1]^3$. The geometric map and
its Jacobian are
\begin{equation}
  \bx(\boldsymbol{\xi}) = \sum_{a} N_a(\boldsymbol{\xi})\, \bx_a ,
  \qquad
  \bJ(\boldsymbol{\xi}) = \frac{\partial \bx}{\partial \boldsymbol{\xi}} ,
  \qquad
  \bA(\boldsymbol{\xi}) = \operatorname{adj} \bJ = \det(\bJ)\, \bJ^{-1} .
  \label{eq:jac}
\end{equation}
Gradients push forward as $\nabla_{\!\bx} \phi = \bJ^{-\mathsf T} \nabla_{\!\boldsymbol\xi}
\phi = \bA^{\mathsf T} \nabla_{\!\boldsymbol\xi} \phi / \det \bJ$.

The sub-control surface associated with parametric direction $d$ is, in the reference
element, a quarter of the unit face normal to $\xi_d$. By Nanson's formula its oriented area
in physical space is
\begin{equation}
  \ba^{(d)} \;=\; \int_{\hat S} \det(\bJ)\, \bJ^{-\mathsf T} \mathbf{e}_d \dd \hat S
  \;=\; \tfrac14 \, \bA_{d,:} \;=\; \tfrac14 \det(\bJ) \, \nabla \xi_d ,
  \label{eq:area}
\end{equation}
the $d$-th row of the adjugate, scaled by the reference sub-face area $\tfrac14$. This single
identity supplies every area vector in the scheme, interior and boundary alike; a boundary
sub-face on the element face normal to $\xi_d$ carries $\pm\ba^{(d)}$ with the sign chosen
outward.

\subsection{Two geometric variants}
\label{sec:geomvariants}

Equation \eqref{eq:area} is evaluated in two ways, and the choice is a modelling decision
rather than an optimisation.

\paragraph{Isoparametric.} $\bJ$, $\bA$ and $\det\bJ$ are evaluated at each sub-control
surface centroid $\boldsymbol{\xi}_s$, and the velocity gradient with them. Each of the twelve
surfaces then carries its own area vector and its own gradient. This is exact for the
trilinear geometry.

\paragraph{Affine.} The geometry is evaluated once, at the element centre
$\boldsymbol{\xi} = (\tfrac12,\tfrac12,\tfrac12)$, and reused for all twelve surfaces. The
four sub-control surfaces of a given direction then share one area vector
$\ba^{(d)}$, and the velocity gradient is a single element constant
\begin{equation}
  \nabla \bu \big|_e \;=\; \frac{1}{\det \bJ} \, \bA^{\mathsf T} \,
  \big[ \, \delta_1 \bu \;\; \delta_2 \bu \;\; \delta_3 \bu \, \big] ,
  \qquad
  \delta_d \bu = \tfrac14 \!\! \sum_{a \in F_d^{+}} \!\! \bu_a \;-\; \tfrac14 \!\! \sum_{a \in F_d^{-}} \!\! \bu_a ,
  \label{eq:affinegrad}
\end{equation}
with $F_d^{\pm}$ the four nodes of the two element faces normal to $\xi_d$. For a
parallelepiped the two variants coincide identically; on a distorted hexahedron the affine
variant is a one-point rule for a bilinear surface and is second-order accurate rather than
exact. Both are carried, because the affine form collapses twelve geometric evaluations into
one and the resulting operator is between three and five times cheaper.

\section{The interior flux}
\label{sec:interior}

Let $s$ be a sub-control surface of an element, separating nodes $i$ and $j$, with area
vector $\ba$ oriented from $i$ to $j$, and write
\begin{equation}
  \bar\bu = \tfrac12 (\bu_i + \bu_j), \qquad
  \bar p  = \tfrac12 (p_i + p_j), \qquad
  \bd = \bx_j - \bx_i .
\end{equation}

\subsection{Mass flux}

The mass flux through $s$ is the linearly interpolated normal flux plus the Rhie--Chow
correction of \S\ref{sec:rc},
\begin{equation}
  \mdot_s \;=\; \rho \, \bar\bu \cdot \ba \;+\; \mdot_s^{\mathrm{RC}} .
  \label{eq:mdot}
\end{equation}

\subsection{Convection, pressure and viscous traction}

The convective momentum flux is upwinded on the sign of $\mdot_s$,
\begin{equation}
  \mdot_s^{+} = \tfrac12\left(\mdot_s + |\mdot_s|_\varepsilon\right), \qquad
  \mdot_s^{-} = \tfrac12\left(\mdot_s - |\mdot_s|_\varepsilon\right),
  \label{eq:mpm}
\end{equation}
so that $\mdot_s^{+} \bu_i + \mdot_s^{-} \bu_j$ selects the upwind node exactly when
$|\cdot|_\varepsilon = |\cdot|$. The pressure is interpolated centrally and the viscous
traction is taken from the element gradient of \S\ref{sec:geomvariants}. The total momentum
and mass fluxes through $s$ are
\begin{align}
  \bF_s &= \mdot_s^{+} \bu_i + \mdot_s^{-} \bu_j + \bar p \, \ba - \bttau \ba ,
  \label{eq:flux} \\
  \bttau \ba &= \mu \left[ \left( \nabla\bu + \nabla\bu^{\mathsf T} \right) \ba \right] .
  \nonumber
\end{align}
Each surface contributes to both of its nodes with opposite signs,
\begin{equation}
  \bR_i \mathrel{+}= \bF_s , \quad
  \bR_j \mathrel{-}= \bF_s , \quad
  R^{c}_i \mathrel{+}= \mdot_s , \quad
  R^{c}_j \mathrel{-}= \mdot_s ,
  \label{eq:scatter}
\end{equation}
which is what makes the scheme discretely conservative: summing $R^c$ over all nodes leaves
only boundary sub-faces.

\subsection{The smoothed upwind switch}
\label{sec:upwind}

The hard switch takes $|\mdot|_\varepsilon = |\mdot|$, which has a corner at $\mdot = 0$.
Where the flow stagnates --- inside a recirculation, or across a partly reversed outlet ---
sub-control surfaces sit arbitrarily close to that corner while the two upwinded states still
differ by $O(1)$, so the flux is not differentiable near the iterate and Newton loses its
rate. Following Venkatakrishnan's requirement that a switch be inactive where it is not
needed and reach that state smoothly, the absolute value is replaced by
\begin{equation}
  |\mdot|_\varepsilon \;=\;
  \begin{cases}
    |\mdot| , & |\mdot| \ge \varepsilon , \\[4pt]
    \dfrac{\mdot^2 \left( 2\varepsilon - |\mdot| \right)}{\varepsilon^2} , & |\mdot| < \varepsilon ,
  \end{cases}
  \label{eq:band}
\end{equation}
which matches $|\mdot|$ in value and slope at $|\mdot| = \varepsilon$, vanishes with zero
slope at the origin, and never exceeds $|\mdot|$. Surfaces carrying real flux are untouched;
surfaces carrying none lose their upwinding and blend to central differencing. The direction
matters: an entropy-fix style band that saturates $|\mdot|$ at a floor $\varepsilon/2$ instead
\emph{adds} dissipation $\tfrac{\varepsilon}{4}(\bu_i - \bu_j)$ on every zero-flux surface,
which in a unidirectional flow is most of the mesh. Taking $\varepsilon = 0$ recovers the hard
switch exactly.

\section{Pressure--velocity coupling}
\label{sec:rc}

On a colocated arrangement the mass flux $\rho\,\bar\bu\cdot\ba$ does not see the odd pressure
mode, and \eqref{eq:cvcont} admits a checkerboard. The Rhie--Chow correction restores the
coupling by adding to the interpolated flux a term proportional to the difference between the
\emph{compact} pressure difference across the surface and the \emph{reconstructed} one,
\begin{equation}
  \mdot_s^{\mathrm{RC}} \;=\; -\,c_s \left[
    \left( p_j - p_i \right)
    - \tfrac12 \left( \nabla p_i + \nabla p_j \right) \cdot \bd
  \right] ,
  \label{eq:rc}
\end{equation}
with the surface coefficient
\begin{equation}
  c_s \;=\; \rho \, D_s \, \frac{|\ba|^2}{\ba \cdot \bd} ,
  \qquad
  D_s \;=\; \frac{\alpha \, |\bd|^2}{2 \mu} ,
  \qquad\text{hence}\qquad
  c_s \;=\; \frac{\alpha \, \rho \, |\bd|^2 \, |\ba|^2}{2 \mu \, (\ba \cdot \bd)} .
  \label{eq:rccoeff}
\end{equation}
Here $\alpha$ is a dimensionless scale, $\alpha = 1$ by default and $\alpha = 0$ recovering the
uncorrected flux. The coefficient is set to zero on a degenerate surface, $|\ba\cdot\bd|$
below a relative floor, so that a collapsed element contributes nothing rather than a
division by a vanishing quantity.

Two properties of \eqref{eq:rc} are worth stating because they determine where it may be
used. First, the bracket annihilates any pressure field whose reconstructed gradient is
exact, in particular any globally linear $p$: the correction is a genuine stabilisation and
not a consistency error. Second, $c_s$ depends only on the geometry and on $\rho, \mu,
\alpha$, so it is constant over a Newton solve and may be precomputed once per surface.

\subsection{The reconstructed nodal gradient}
\label{sec:nodalgrad}

The gradients in \eqref{eq:rc} are the volume-weighted average of the element gradients
meeting at the node,
\begin{equation}
  \nabla p_i \;=\;
  \frac{\displaystyle\sum_{e \,\ni\, i} |\det \bJ_e| \; \nabla p \big|_e}
       {\displaystyle\sum_{e \,\ni\, i} |\det \bJ_e|} ,
  \label{eq:nodalgrad}
\end{equation}
where $\nabla p|_e$ is the element-constant gradient \eqref{eq:affinegrad} applied to $p$, or
its isoparametric counterpart at the element centre. This is a linear operator
$\mathcal{G} : \mathbb{R}^{N} \to \mathbb{R}^{3N}$ acting on the nodal pressures, and its
stencil reaches every node of every element containing $i$ --- a wider set than the
nearest-neighbour graph of the mesh. That single fact separates two of the operators in
\S\ref{sec:operators}.

\section{Boundary closure}
\label{sec:boundary}

A control volume touching $\partial\Omega$ is not closed by interior sub-control surfaces
alone; its boundary sub-faces must carry the remaining flux, or \eqref{eq:cvcont} is not a
balance at all. For a boundary sub-face of node $i$ with outward area vector $\ba$ the flux is
$\mdot = \rho\, \bu_i \cdot \ba$ and three closures are used.

\paragraph{Closed / prescribed velocity.} The full flux with the local state,
\begin{equation}
  \bR_i \mathrel{+}= \mdot\, \bu_i + p_i \ba - \bttau \ba ,
  \qquad
  R^c_i \mathrel{+}= \mdot .
  \label{eq:bstd}
\end{equation}
On a no-slip wall $\bu_i = 0$ makes the convective and mass terms vanish identically and
leaves the pressure and viscous traction, which is the correct wall contribution.

\paragraph{Prescribed pressure.} As \eqref{eq:bstd} with $\bar p$ in place of $p_i$. The
viscous traction is still evaluated from the interior state, so this prescribes the pressure
and not the whole normal traction.

\paragraph{Natural (do-nothing) outflow.} The condition $(p\,\bI - \bttau)\bn = \bt$ is
imposed by \emph{omitting} the pressure and viscous terms rather than evaluating them,
\begin{equation}
  \bR_i \mathrel{+}= \max(\mdot, 0)\, \bu_i + \bt \, |\ba| ,
  \qquad
  R^c_i \mathrel{+}= \mdot .
  \label{eq:bnat}
\end{equation}
Two details are load-bearing. Retaining $p_i \ba$ here would leave the constant-pressure mode
in the null space --- a uniform pressure shift integrates to zero over every closed control
volume --- so the gauge would remain undetermined and require a pin. And the convective term
is clipped to $\max(\mdot,0)$ because this closure has no upwind switch: on a sub-face with
$\mdot < 0$ the unguarded form convects the exterior state inward, which is the classical
backflow instability of a recirculating outlet. The continuity row keeps the true, unclipped
$\mdot$, since clipping it would destroy global mass conservation.

\section{Time integration and body force}

The transient term is the only place where the control-volume structure makes the mass matrix
trivial. In \eqref{eq:cvmom} the volume integral of $\rho\bu$ over $C_i$ is $\rho V_i \bu_i$
up to the same order as the rest of the scheme, so the mass matrix \emph{is} the control
volume: diagonal, already available from \eqref{eq:nodevol}, and requiring no assembly. A
consistent finite-element mass matrix would be a different discretisation, not a better one.

With backward differentiation of order $k$ and step $\Delta t$,
\begin{equation}
  \bR_i \mathrel{+}= \frac{\rho V_i}{\Delta t}
  \left( a_0 \bu_i^{n+1} + a_1 \bu_i^{n} + a_2 \bu_i^{n-1} \right),
  \qquad
  \begin{array}{lccc}
    \text{BDF1:} & a_0 = 1, & a_1 = -1, & a_2 = 0, \\
    \text{BDF2:} & a_0 = \tfrac32, & a_1 = -2, & a_2 = \tfrac12,
  \end{array}
  \label{eq:bdf}
\end{equation}
with $\sum_k a_k = 0$, so a constant history contributes exactly nothing. The term touches the
three velocity components and never the continuity row: the incompressibility constraint has
no time derivative, and giving it one would be a different system. Its contribution to the
Jacobian is therefore the diagonal
\begin{equation}
  \frac{\partial}{\partial \bu_i}
  \left[ \frac{\rho V_i a_0}{\Delta t} \bu_i \right]
  = \frac{\rho V_i a_0}{\Delta t} \, \bI_{3\times3} ,
  \label{eq:bdfjac}
\end{equation}
independent of the history. A body force enters the same way, as $\bR_i \mathrel{-}=
\rho\,\bff(\bx_i)\, V_i$, and for the same reason: neither term involves a sub-control-surface
flux, so both are node-local post-passes over the mesh.

\section{The discrete operators}
\label{sec:operators}

Collect the unknowns as $w = (\bu_1,p_1,\dots,\bu_N,p_N) \in \mathbb{R}^{4N}$. Everything
above defines one nonlinear map; what a solver actually evaluates is the family below. They
are listed with what each carries, because they do not all carry the same terms and the
differences are deliberate.

\subsection{$\mathcal{F}$ --- the residual}

\begin{equation}
  \mathcal{F}(w)_i \;=\;
  \underbrace{\sum_{s \,\in\, \partial C_i} \pm \bF_s}_{\text{interior sub-control surfaces}}
  \;+\; \underbrace{\sum_{b} \bF_b}_{\text{boundary closure}}
  \;+\; \underbrace{\frac{\rho V_i}{\Delta t}\sum_k a_k \bu_i^{n+1-k}}_{\text{transient}}
  \;-\; \underbrace{\rho \bff_i V_i}_{\text{body force}} ,
  \label{eq:residual}
\end{equation}
with the continuity component assembled from $\pm\mdot_s$ and $\mdot_b$ in the same sweep.
Evaluating it requires the nodal gradient \eqref{eq:nodalgrad} of the current pressure, which
is a function of the Newton iterate and may therefore be computed once per Newton step rather
than once per residual.

\subsection{$\mathcal{J}(w)\,v$ --- the exact Jacobian action}
\label{sec:jacaction}

The Fr\'echet derivative applied to a direction $v = (\bv, q)$. Every term of
\eqref{eq:residual} is differentiated, including the reconstruction \eqref{eq:nodalgrad}:
since $\mathcal{G}$ is linear, its derivative in the direction $q$ is $\mathcal{G} q$, and the
Rhie--Chow perturbation is
\begin{equation}
  \delta \mdot_s^{\mathrm{RC}} \;=\; -\,c_s \left[
    \left( q_j - q_i \right)
    - \tfrac12 \left( \nabla q_i + \nabla q_j \right) \cdot \bd
  \right] ,
  \qquad \nabla q = \mathcal{G} q .
  \label{eq:rcjac}
\end{equation}
The upwind switch is differentiated through \eqref{eq:band}, whose derivative is
\begin{equation}
  \frac{\dd |\mdot|_\varepsilon}{\dd \mdot} =
  \begin{cases}
    \operatorname{sgn}(\mdot) , & |\mdot| \ge \varepsilon , \\[4pt]
    \operatorname{sgn}(\mdot)\,\dfrac{4\varepsilon |\mdot| - 3 \mdot^2}{\varepsilon^2} , & |\mdot| < \varepsilon .
  \end{cases}
\end{equation}
This operator is the one a Krylov method applies, and \eqref{eq:rcjac} is what makes it
expensive: $\mathcal{G}q$ is a full sweep over the elements, the direction changes at every
iteration, and so --- unlike $\mathcal{G}p$ --- it cannot be hoisted out of the solve.

\subsection{$J(w)$ --- the assembled Jacobian, with frozen reconstruction}
\label{sec:jacassembled}

The same derivative with the second term of \eqref{eq:rcjac} dropped,
\begin{equation}
  \delta \mdot_s^{\mathrm{RC}} \;=\; -\,c_s \left( q_j - q_i \right) ,
  \label{eq:rcfrozen}
\end{equation}
that is, the reconstruction is held frozen at the current iterate. This is not an oversight
and not an approximation made for speed: by \S\ref{sec:nodalgrad} the operator $\mathcal{G}$
couples node $i$ to every node of every element containing $i$, so retaining it would widen
the sparsity pattern beyond the mesh's nearest-neighbour graph. The assembled matrix exists to
build preconditioners, and a preconditioner on a wider graph than the operator it
preconditions is a poor trade. With \eqref{eq:rcfrozen} the pattern is exactly the
node-to-node graph and the matrix has the $4\times4$ block structure
\begin{equation}
  J \;=\;
  \begin{pmatrix}
    F & B^{\mathsf T} \\
    B & C
  \end{pmatrix} ,
  \label{eq:blockform}
\end{equation}
with $F$ the momentum block (convection, viscosity and, when unsteady, the mass term), $B$ the
discrete divergence and $B^{\mathsf T}$ the discrete gradient. The pressure block carries a
plus sign here because the continuity residual is written as $+\oint \rho\,\bu\cdot\bn$; the
$-C$ familiar from the stabilised-Stokes literature is the same operator with that row
negated.

$C$ is the Rhie--Chow operator, and differentiating \eqref{eq:rcfrozen} through the scatter
\eqref{eq:scatter} identifies it exactly: $C_{ij} = -c_s$ for each sub-control surface joining
$i$ and $j$, and $C_{ii} = \sum_{s \ni i} c_s$. It is the graph Laplacian of the mesh with
weights $c_s > 0$ --- symmetric, positive semidefinite, with the constant vector in its
kernel, which is the discrete statement that the pressure is determined only up to a gauge.

Without Rhie--Chow the fourth row and column of $J$ do not merely lose their diagonal: they
are structurally empty at every node whose control volume is closed. Summing the oriented
areas over a closed surface gives $\sum_{s \in \partial C_i} \sigma_s \ba_s = 0$, and both
$B_{ii}$ and $B^{\mathsf T}_{ii}$ are $\tfrac{\rho}{2}$ and $\tfrac12$ times exactly that sum.
So $C = 0$ leaves a genuine saddle point whose $4\times4$ diagonal blocks are singular with a
one-dimensional kernel --- which is why no point- or block-wise relaxation can be built on
them, and why the boundary closure of \S\ref{sec:boundary} is a prerequisite for the
preconditioner and not only for conservation.

\subsection{$D$ --- the block diagonal}

The $N$ diagonal $4\times4$ blocks of \eqref{eq:blockform},
\begin{equation}
  D_i \;=\; \frac{\partial \mathcal{F}_i}{\partial w_i} \;=\;
  \begin{pmatrix}
    F_{ii} & B^{\mathsf T}_{ii} \\
    B_{ii} & C_{ii}
  \end{pmatrix}
  \in \mathbb{R}^{4\times4} ,
\end{equation}
carrying the same terms as $J$, including the boundary closure and the transient diagonal
\eqref{eq:bdfjac}. It is the object a block-Jacobi smoother inverts, and by the closure
identity of \S\ref{sec:jacassembled} its usefulness rests entirely on $C_{ii} \neq 0$: on a
closed control volume the other three entries of its fourth row and column vanish, so without
the Rhie--Chow term $D_i$ is singular for every $i$ and a smoother built on it is undefined
rather than merely slow.

\subsection{$J v$ --- the sparse matrix--vector product}

Given $J$ assembled, the action can also be evaluated as a block sparse matrix--vector
product. Mathematically this is \S\ref{sec:jacassembled} applied to $v$ --- it therefore
carries the frozen reconstruction \eqref{eq:rcfrozen} and differs from \S\ref{sec:jacaction}
by exactly the term $\tfrac12 c_s (\nabla q_i + \nabla q_j)\cdot\bd$. Its cost is set by the
sparsity pattern alone and is independent of which physical terms the entries carry.

\subsection{$S$ --- the pressure Schur complement diagonal}

For a SIMPLE-type preconditioner the exact Schur complement $C - B F^{-1} B^{\mathsf T}$ is
replaced by its diagonal, with $F$ approximated by its own velocity diagonal,
\begin{equation}
  S_i \;=\; \operatorname{diag}\!\left(
    C - B \, \operatorname{diag}(F)^{-1} B^{\mathsf T}
  \right)_i ,
  \label{eq:schur}
\end{equation}
read directly from the assembled blocks: $C_{ii}$ is the pressure--pressure entry of block
$(i,i)$, and the correction sums over the row of $J$. Unlike the classical $\mu\,M_p^{-1}$
approximation this makes no assumption about which term dominates --- it reads both the
Rhie--Chow operator and the velocity coupling out of the matrix that was assembled.

\subsection{Splitting the momentum block}

The viscous contribution to \eqref{eq:flux} is $-\bttau\ba$ with $\bttau$ linear in $\bu$ and
$\ba$ purely geometric, so that part of $F$ depends on the mesh and $\mu$ alone and is
constant across a Newton solve. Writing $F = F_{\mathrm{visc}} + F_{\mathrm{conv}}(\bu)$
allows the constant half to be formed once and the varying half added at each iteration; the
sum reproduces $F$ exactly, since the two are the two halves of one expression rather than two
approximations. The Rhie--Chow blocks belong entirely to the varying half.

\section{What the terms cost}
\label{sec:cost}

The operators above differ in cost by more than an order of magnitude, and the differences do
not follow the number of terms. The table gives measured throughput in millions of degrees of
freedom per second on one 72-core Grace socket at $4.12 \times 10^6$ degrees of freedom, with
the bare element sweep --- no Rhie--Chow, no boundary closure, no transient --- as the
reference.

\begin{center}
\begin{tabular}{lrr}
\toprule
Operator & MDOF/s & vs.\ bare sweep \\
\midrule
Residual $\mathcal{F}$, interior flux only            & 2945 & 1.00 \\
\quad $+$ Rhie--Chow, reconstruction reused           & 1759 & 1.67 \\
\quad $+$ boundary closure                            & 1559 & 1.89 \\
\quad $+$ transient                                   & 1436 & 2.05 \\
\quad $+$ reconstruction rebuilt every evaluation     &  667 & 4.41 \\
\midrule
Jacobian action $\mathcal{J}v$, interior flux only    & 2067 & 1.00 \\
\quad $+$ exact Rhie--Chow \eqref{eq:rcjac}           &  605 & 3.42 \\
\quad $+$ boundary closure and transient              &  561 & 3.68 \\
\midrule
Matrix--vector product $Jv$                           &  462 & 4.47 \\
Assembly of $J$                                       & 39--87 & 34--75 \\
Block diagonal $D$                                    &   90 & --- \\
\bottomrule
\end{tabular}
\end{center}

Three readings of this table bear on the design rather than on the implementation.

The exact Rhie--Chow term costs the Jacobian action a factor of $3.4$ on its own, against
$1.67$ for the same term in the residual. The asymmetry is \eqref{eq:rcjac}: the residual
reconstructs the gradient of a field that is fixed over the Newton step, the action
reconstructs the gradient of a direction that changes at every Krylov iteration, and only the
former can be amortised. Where the reconstruction cannot be reused --- the last residual row
--- the residual pays the same factor.

The matrix--vector product is slower than the matrix-free action carrying the same physics,
by $1.25$, before counting the assembly that produced the matrix. The reason is traffic
rather than arithmetic: the assembled blocks are roughly $840$ bytes per degree of freedom
and must be streamed in full at every application, whereas the matrix-free operator reads the
state and the mesh and recomputes the rest. This is also why the product's cost is flat in the
physics while the matrix-free operator's is not. The assembled matrix earns its place as the
source of $D$ and $S$, not as a way to apply $J$.

Adding physics narrows the gap between matrix-free and assembled from $4.5$ to $1.25$ without
closing it. Any future term heavy enough to close it would also, by the same argument, be
heavy enough to make assembly worth revisiting --- which is the quantitative form of the
usual matrix-free trade-off, and the reason the terms are enumerated separately here rather
than folded into a single residual.

\end{document}
